% Scope: Explain static, RF, gradient, acoustic and cryogenic hazards from their mechanisms; keep operational limits tied to current authoritative standards when drafted.
% Status: architecture only; no chapter prose or completed problems yet.
% Prerequisites: Chapters 1, 17, 18 and 24.
% Planned worked examples: RF duty-cycle scaling and gradient slew comparisons without patient-specific clearance.
% Planned figures: Energy deposition and induced electric-field pathways.
% Planned challenge problems: Compare equal-flip pulses with unequal heating; distinguish SAR from local temperature.
\chapter{Safety and System Limits}
\label{ch:26}

\section{Static-field interactions}
\label{sec:26:static-field-interactions}
\subsection{Forces, torques and projectile mechanisms}
\subsection{Spatial gradients and active shielding}
\subsection{Implant interactions and screening principles}

\section{RF energy deposition}
\label{sec:26:rf-energy-deposition}
\subsection{Electric fields, dissipated power and SAR}
\subsection{Global versus local heating and duty cycle}
\subsection{Conductive loops, implants and model uncertainty}

\section{Gradient and acoustic effects}
\label{sec:26:gradient-and-acoustic-effects}
\subsection{Induced electric fields and PNS}
\subsection{Waveform dependence beyond peak slew}
\subsection{Mechanical forces, acoustic noise and protection}

\section{Cryogenics and system failures}
\label{sec:26:cryogenics-and-system-failures}
\subsection{Stored energy and quench mechanisms}
\subsection{Cryogen release and oxygen displacement}
\subsection{Fault response and system interlocks}

\section{Limits and responsible interpretation}
\label{sec:26:limits-and-responsible-interpretation}
\subsection{Scanner operating modes and standards}
\subsection{MR Safe, MR Conditional and MR Unsafe labeling}
\subsection{Physics estimates versus device-specific conditions}

\section{Worked examples}
\label{sec:26:worked-examples}
% Planned calculations: RF duty-cycle scaling and gradient slew comparisons without patient-specific clearance.
\subsection{Derivation and quantitative calculation}
\subsection{Assumptions, limiting cases and scanner interpretation}

\section{Challenge problems}
\label{sec:26:challenge-problems}
% Problem design: Compare equal-flip pulses with unequal heating; distinguish SAR from local temperature.
% Use challengeproblem and stable labels prob:26:<topic>.
% Complete solutions belong in Appendix E, section for this chapter.
% Use \begin{solution}{prob:26:<topic>} ... \end{solution} there.
\subsection{Derivation and quantitative design}
\subsection{Model criticism and integrated reasoning}
